\begin{flushright}
  \fechaclase{05 de noviembre de 2026}
\end{flushright}

\subsection{SMC con convergencia en tiempo prescrito}

Sea el sistema SISO

\begin{equation*}
    \begin{aligned}
        \dot{x}_1 &= x_2 \\
        \dot{x}_2 &= f(x, t)  u\\
        y&= x_1
    \end{aligned}
\end{equation*}

Se define el error de seguimiento y sus derivadas como

\begin{equation*}
    \begin{aligned}
        e_1 &= r-y = r-x_1 \\
        e_2 &= \dot{r} -x_2 \\
        \dot{e}_1 &= \dot{r}- \dot{x}_1 = \dot{r}- x_2 = e_2 \\
        e_2 &= \ddot{r}- \dot{x}_2 = \ddot{r}-f(x,t)-u
    \end{aligned}
\end{equation*}

el objetivo es que $e_1=0$ y $e_2=0$ cuando $t\leq T$, para algún tiempo T arbitario. Se introduce la funcion de error desaeado

\begin{equation*}
    \begin{aligned}
        e_1^*(t)&\\
        e_2^*(t)&=\dot{e}_1^*(t)
    \end{aligned}
\end{equation*}

estas funciones deben cumplir

\begin{equation*}
    \begin{aligned}
        e_1^*(0) &= a \\
        e_2^*(0) &= b \\
        e_1^*(T) &= 0 \\
        e_2^*(T) &= 0 \\
        \dot{e}_1^*(0) &= b 
    \end{aligned}
\end{equation*}

y para $t \geq 0$ $e_1^*(t)=0$ y $e_2^*(t)=0$. Se busca un polinomio para $e_1^*(t)$ donde $0 \leq t \leq T$ cumpla con $\tau = t/T\epsilon [0,1]$ bajo las condiciones

\begin{enumerate}
    \item raiz doble \\
        \begin{equation*}
            e_1^* = (1-\tau)^2(\alpha+\beta \tau)
    \end{equation*}

    \item como $e_1^*(0)=a$, entonces en $\tau=0:e_1^*=\alpha=a$
    
    \item como $\dot{e}_1^*(0)=b$, entonces en $\frac{d}{dt}\tau=\frac{1}{T}$ y $\dot{e}_1^*=\frac{1}{T}[-2(1-\tau)(\alpha-\beta \tau)+\beta(1-\tau)^2]$ en $\tau=0:\dot{e}_1^*=\frac{1}{T}(-2a+\beta)=b$ para que $\beta=bT+2a$
        
\end{enumerate}
 